\documentclass{article}
\PassOptionsToPackage{numbers,sort&compress}{natbib}
\usepackage[preprint]{neurips_2026}
\makeatletter
\renewcommand{\@noticestring}{Anonymous research manuscript.}
\makeatother
\usepackage[T1]{fontenc}
\usepackage{amsmath,amssymb,booktabs,microtype}
\usepackage[hidelinks]{hyperref}

\title{MinPI for Feature Selection with a Limited Number of Refits}
\author{Anonymous Author(s)}
\begin{document}
\maketitle
\begin{abstract}
Feature importance gives an order in which to try columns. Choosing how many to retain can still require one model fit per column. We study a feature selector that ranks columns by which thresholds they can represent for one another. A column earns a vote when adding it to a coalition covers a previously unrepresented threshold. MinPI takes the smallest of its relative voting shares at three coalition sizes. The votes have an exact counting formula, and cumulative covered weight identifies at most six widths for validation. On four held-out datasets, this selector took less time than refitting every prefix of an XGBoost gain ranking. It also kept a comparable or smaller number of columns on three datasets. Two datasets met the joint criterion against Gain with a bounded search. Replacing the minimum with a mean gave identical selected columns on two datasets across all five seeds. Much of the measured speed difference came from testing fewer widths.
\end{abstract}

\section{Introduction}
A tree split assigns its training gain to a single input even when another input would make almost the same split. This assignment matters when selecting a few columns from many related measurements. The rank of an individual column gives little indication of how much it can add once its substitutes have been selected. Checking the prediction error after each change can settle that question on validation data, but on a path of $p$ widths it can also require $p$ separate training fits.

We adapt the voting and minimum operations of the min-marginal power index~\citep{shi2026minpi} to select columns for retraining. The voters are threshold partitions of training observations. A candidate casts a marginal contribution to a coalition's covered weight whenever it can represent a partition that the coalition misses. We compute the expected marginal contribution at each of three coalition sizes and choose the column with the largest worst-size share. Once selected, a column's partitions leave the game. The resulting order and cumulative coverage determine a small number of models to train.

Game-theoretic feature selection has previously used model predictions on representative samples to estimate feature contributions~\citep{zhaoStrategicInputs}. Here we use ordinary tree gain and split agreement to define a simpler game in which the coalition average has a closed form. Its utility measures coverage of an empirical collection of thresholds. Whether a downstream learner trained on those columns predicts well is an experimental question. We compare actual refits against both the exhaustive path and inexpensive searches using the same fitted model's gain ranking. This second comparison tests how much the votes add once the cost of searching every width has been removed.

\section{Split coverage votes}

Consider a training matrix with $p$ columns. Two financial ratios, for example, can place almost the same firms above and below a threshold. Fitting a tree on one ratio can give that ratio the full gain from the split. To measure which ratios can substitute for one another, we compare threshold partitions constructed on a common sample of training rows.

\paragraph{Constructing the voters.}
Fit XGBoost on the training rows and keep the total split gain $G_j$ attributed to each column~$j$~\citep{chen2016xgboost}. On a uniform sample of at most 512 training rows, take the seven empirical octiles of each column's finite values. A threshold indicator equals one for an observed value at or below its cut. Add the missingness indicator if applicable and discard constant or repeated indicators within a column. Each retained indicator $b$ is a voter. A column represents $b$ if one of its indicators, in either orientation, makes at most half as many row disagreements with $b$ as the better constant prediction. Write $R_b$ for its representatives. This compares splits on all sampled rows. A tree surrogate split usually measures agreement among the observations reaching one node~\citep{seifert2019surrogate}.

Give each voter a weight $w_b$ by dividing its owner's $G_j$ among that owner's indicators in proportion to their between-group sums of squares for the target. Use one-hot targets for classification, summing over classes. Split the gain equally if all sums vanish. For a set $S$ of selected columns the covered weight is
\begin{equation}
 F(S)=\sum_b w_b\,\mathbf{1}\{S\cap R_b\ne\varnothing\}\,.
 \label{eq:coverage}
\end{equation}
This is a weighted coverage game. Feature selection by coverage and facility location predates this work~\citep{liu2013coverage}. Here each threshold votes for the columns that represent its partition. Computing all pairs of agreements for $B$ thresholds on $m$ sampled rows costs $O(mB^2)$ operations and the implementation stores a $B$-by-$B$ matrix. This cost can dominate at high input dimension.

\paragraph{Choosing a column.}
Remove the voters covered by the current selection and let $n$ be the number of remaining columns. Fix a candidate $j$ and draw a set of $k$ from the other $n-1$ columns uniformly. Voter $b$ needs $j$ precisely when $j\in R_b$ and the draw misses its other representatives. The expected weight gained by adding $j$ and its share of total marginal weight at size $k$ are
\begin{equation}
 c_j(k)=\sum_{b:\,j\in R_b}w_b
        \frac{\binom{n-|R_b|}{k}}{\binom{n-1}{k}},
 \qquad
 q_j(k)=\frac{c_j(k)}{\sum_{\ell}c_\ell(k)}\,.
 \label{eq:pivotal}
\end{equation}
Indeed, the numerator of the binomial ratio counts the draws using only columns outside $R_b$. This evaluates the coalition average exactly for the sampled voters. No further model fit is required.

Use the distinct integers nearest $(n-1)/6$, $(n-1)/2$ and $5(n-1)/6$, rounding half-integers to even and omitting sizes with zero total marginal weight. Append the column with the largest $\min_k q_j(k)$ and repeat. Higher average tree split gain, then higher column index, breaks ties and completes the order once the remaining votes vanish. The construction retains the minimum over coalition contexts of the min-marginal approach~\citep{shi2026minpi}. Its votes concern the extra split weight represented by including a column.

Normalization by size matters. Each raw $c_j(k)$ decreases weakly with $k$, so its minimum always occurs at the largest tested size. If all remaining positive-weight voters have the same number of representatives, the common binomial ratio cancels and each $q_j(k)$ is constant across sizes. With unequal multiplicities the minimum and mean can choose different columns. Give $A$ a private vote of weight 5, $B$ and $C$ a shared vote of weight 11, and $D$ and $E$ private votes of weight 1. Across $k=1,2,3$ the smallest share of $A$ is $10/47$, just below the $11/50$ of $B$ and $C$. The mean ranks $A$ first. This calculation establishes a difference in ranking. Measuring its effect on predictions requires refitting on the selected columns.

\paragraph{Choosing a width.}
For each prefix of length $h$, compute $F(S_h)$ from the voters first covered along the order. Fit the first prefix reaching each of $1/2,3/4,7/8,15/16$ of total weight and reuse the full training fit at width $p$. For zero total mass use the corresponding fractions of $p$. Find the smallest tested width within 0.005 of the best validation utility. Try one intermediate width halfway to the previous tested width (or to zero) if the gap exceeds one, then apply the same selection rule to all scores. A nonmonotone utility curve can attain its optimum between tested widths. The procedure evaluates at most $\min(p,6)$ distinct widths and fits the selected support once more on training and validation rows for deployment. The same fit bound holds for an ordinary Gain ranking with the same search schedule.

\section{Evaluation}

Each method ranks the same input matrix, chooses a prefix on validation data, and refits a predictor using only its selected columns. All fits use a depth-four XGBoost model with 100 boosting rounds, learning rate 0.1, native missing values and one CPU thread on the same CUDA device. We charge the initial fit, vote construction, search, final refit and test prediction to each method independently. Timing excludes common input loading and writing predictions to disk. The run order rotates across seeds.

The direct comparator orders columns by average training split gain, refits every prefix and selects the narrowest validation maximizer. We also record the smallest width within 0.005 of its maximum. Three additional controls isolate the components of the search. Pooled voting replaces $\min_k q_j(k)$ with its arithmetic mean and uses the identical refit policy. Mass-scheduled Gain uses the ordinary Gain order, assigning each column its own total split gain for width selection. Geometric Gain tries powers of two and the full set, followed by at most two bisections below the observed near-best width. All three controls use the same 0.005 allowance.

The timing and width criteria apply to each dataset. On each seed we record the time difference and the excess width after allowing the larger of one column and five percent of the comparator's width. A dataset meets the operational target if the median paired time difference is negative and the median adjusted width excess is nonpositive. We require this on at least 80\% of datasets in a panel. We separately count individual seeds satisfying both conditions, and datasets on which the median primary score falls by at most~0.01. The primary score is binary or macro one-vs-rest AUC for classification and $R^2$ for regression. We report accuracy and MAE in their native units. Validation truncates negative $R^2$ at zero in every arm, following the original selector. Test-set $R^2$ is untruncated.

\paragraph{Data.}
The six development datasets are Abalone, Credit Card Fraud, Gisette, ISOLET, Madelon and Spambase. All use the existing prepared columns and partitions. We then fixed four UCI collections and their splits before fitting a model to them~\citep{uciMoCap,uciBankruptcy,uciTurbine,uciTemperature}. MoCap Hand Postures has 36 coordinates and 2,800 frames, with at most 200 frames per person and disjoint subjects. Taiwanese Bankruptcy has 95 ratios and 6,819 financial records. Identical input vectors travel together across splits. The released table lacks firm and time identifiers. Turbine CO prediction uses nine plant and ambient readings, retaining every tenth observation in each year (3,676 rows). Seoul temperature forecasting uses 21 inputs and 7,723 station-days, with both observed future extrema removed. The last year is held out in each regression task. Both regressions reuse that same test year across five seeds, so these ten scores describe two test periods. The turbine release concerns another endpoint at a plant represented during earlier development.

For the two classifications, models fit on approximately 48\% of groups, choose widths on the next 32\% and score the remaining 20\%. On the time splits the first two years fit the learner, the next two years jointly select widths, and refitting uses all four development years. Every arm sees the same rows. Fit-derived choices and forecast files were recorded before test scoring.

\paragraph{External results for the split-coverage selector.}
The selector ran faster than exhaustive Gain on all four collections. Three met the joint width and time criterion (Table~\ref{tab:base-external}). Seoul temperature prediction retained two more columns for a median $R^2$ gain of 0.0024, exceeding the one-column allowance. The registered threshold required all four datasets. Across individual seeds, 16 of 20 met both operational criteria.

\begin{table}[htbp]
\centering
\caption{Five-seed medians for split coverage (MinPI) and exhaustive Gain. Each pair lists MinPI first. Primary scores are test AUC for the first two rows and test $R^2$ for the last two. The companion scores are accuracy and MAE, respectively. Counts give the number of seeds simultaneously meeting the width and timing criteria.}
\label{tab:base-external}
\small
\begin{tabular}{lrrrrr}
\toprule
Dataset & Columns & Seconds & Primary score & Companion score & Count\\
\midrule
Hand postures & 28 / 29 & 4.70 / 24.23 & .925 / .927 & .703 / .705 & 5/5\\
Bankruptcy    & 30 / 61 & 1.38 / 17.19 & .941 / .947 & .975 / .971 & 4/5\\
Turbine CO    & 2 / 2   & .77 / 1.48 & .576 / .576 & .965 / .965 & 5/5\\
Temperature   & 16 / 14 & 1.07 / 3.37 & .657 / .655 & 1.362 / 1.388 & 2/5\\
\bottomrule
\end{tabular}
\end{table}

The comparator's width rule explains part of the apparent reduction. On Bankruptcy, exhaustive Gain retained a median of 61 columns at its validation maximum and 27 with the same 0.005 allowance as MinPI. Only hand postures and turbine CO still passed the width criterion after this change. Against each practical control (pooled voting, mass-scheduled Gain and geometric Gain), two of four datasets passed the joint criterion. On hand postures and temperature the minimum and mean deployed exactly the same columns and made the same predictions in all ten paired runs. On turbine CO the minimum improved $R^2$ by a paired median of 0.050 over the mean. Mass-scheduled Gain selected the very same two columns as the minimum on every seed. These comparisons identify a large saving from restricting the number of retrained widths. They supply little evidence that taking the smallest vote accounts for that saving.


\bibliographystyle{plainnat}
\bibliography{references,split-vote}
\end{document}
